\scanpage{121}

\handwritten{4.} Revenons à la situation de 2, et supposons que la structure de $(\handwritten{\bar}{X},\handwritten{\bar}{U})$ soit celle envisagée dans 3, et que de même $\handwritten{\bar}{Y}$ soit un schéma local régulier et $\handwritten{\bar}{V}=\handwritten{\bar}{Y}-\operatorname{supp}\handwritten{\bar}{E}$, $\handwritten{\bar}{E}$ étant un diviseur à croisements normaux sur $\handwritten{\bar}{Y}$. Soit $C'$ l'ensemble des composantes irréductibles de $\handwritten{\bar}{E}$, On voit alors que pour tout $c'\handwritten{\in}C'$, l'image inverse de $\handwritten{\bar}{E}_{c'}$ est un diviseur sur $\handwritten{\bar}{X}$ dont le support est contenu dans celui de $\handwritten{\bar}{\unsure{\deleted{E}D|\deleted{D}E}}$, donc est de la forme
\[
\handwritten{\bar}{f}(\handwritten{\bar}{E}_{c'})=\handwritten{\sum}_{c}\unsure{a|n}_{c,c'}\handwritten{\bar}{D}_c\ ,
\]
où pour $c'$ fixé, les $n_{c,c'}$ sont des entiers $\handwritten{\geq}0$ non tous nuls. Ces notations fixées, un calcul immédiat \deleted{\unsure{de rev|de re}} d'images inverses de revêtements kummériens montre que pour deux points correspondants $\handwritten{\xi}$ et $\handwritten{\eta}$, \deleted{l'homom} on a commutativité dans le diagramme
\diagram{\(\handwritten{(4.1)}\qquad\)
\begin{tikzcd}[ampersand replacement=\&,column sep=large,row sep=large]
\&\handwritten{\deleted{\unsure{H_\xi|G_\xi}}}\&G_{\handwritten{\xi}}^{\handwritten{\deleted{\unsure{\xi|x}}}}\handwritten{\arrow}[r,"{\handwritten{\phi}}","{\deleted{\unsure{N|\phi}}}"']\&H_{\handwritten{\eta}}\\
\handwritten{-}\&{}\handwritten{\arrow}[u,gray]\&\handwritten{\prod}_{\handwritten{\ell}}\underline{Z}_{\handwritten{\ell}}(1)(\handwritten{\bar}{x})^C\handwritten{\arrow}[u]\handwritten{\arrow}[r,"N"]\&\handwritten{\prod}^{\handwritten{\prime\prime}}_{\handwritten{\ell}}\underline{Z}_{\handwritten{\ell}}(1)(\handwritten{\bar}{y})^{C'}\handwritten{\arrow}[u]
\end{tikzcd}\(\ ,\)}
les flèches \deleted{horizontales}\inserted{verticales} étant les épimorphismes canoniques, et $N$ étant l'homomorphisme défini par la matrice $N=(n_{c,c'}\handwritten{)}_{(c,c')\handwritten{\in}C\mathbin{\times}C'}$, compte tenu que l'homomorphisme $k(\handwritten{\bar}{y})\handwritten{\longrightarrow}k(\handwritten{\bar}{x})$ induit des isomorphismes
\[
\underline{Z}_{\handwritten{\ell}}(1)(\handwritten{\bar}{y})\handwritten{\xrightarrow{\sim}}\underline{Z}_{\handwritten{\ell}}(1)(\handwritten{\bar}{x})\ .
\]
Pour que \inserted{l'homomorphisme} $N$ soit un isomorphisme, il est donc nécessaire et suffisant que \deleted{det $N$} \inserted{$\operatorname{card}C=\operatorname{card}C'$ et que} l'entier ordinaire $\text{dét }N$ soit égal à $\pm p^r$, où $p$ est l'exposant caractéristique de $k(x)$ et $k(y)$; dans ce cas, on conclut donc que les \inserted{sous-}groupes d'inertie de $H$ en $y$ sont les conjugués des images des \inserted{sous-}groupes d'inertie de $G$ en $x$.

\margin{\handwritten{$\unsure{\bar C|\bar E}=\unsure{Y^u|Y^{\vee}|\check Y}$}}
5. Supposons en particulier que l'on ait $V=Y-\operatorname{supp}E$, $E$ étant un diviseur \inserted{$\handwritten{\geq}0$} sur $Y$, \deleted{tel que} $U=f^{\handwritten{-1}}(V)$, et que \deleted{pour tout $\unsure{x\handwritten{\in}|x=}f^{\handwritten{-1}}(E)$, posant} les conditions suivantes soient satisfaites:

\inserted{\handwritten{a) Pour $x\in f^{-1}(\operatorname{supp}E)$, $y=f(x)$, $\underline{O}_{X,x}$ et $\underline{O}_{Y,y}$ sont \unsure{noethériens|noethiriens}.}}

\deleted{\unsure{a}}\handwritten{b}) Le diviseur \inserted{\handwritten{$D=$}} $f^{\handwritten{*}}(E)$ est défini, c'est un sous-schéma \uline{régulier} de
% archive sheet number: 120
